% =============================================================================
% Section 8: Conclusion
% =============================================================================
\section{Conclusion}\label{sec:conclusion}

The International Sovereignty Index applied to the EU-27 produces a
ranking that is structurally determined by two of its six component
axes.  Defence and logistics together account for 69.4\% of
cross-country composite variance.  In every member state, one of
these two axes is the highest-scoring component.  Removing the
defence axis produces the largest rank disruption of any single-axis
omission.  The remaining four axes---financial, energy, technology,
and critical inputs---collectively account for 30.6\% of variance
and do not dominate any country's profile.

This is not a design flaw.  It is an empirical property of the
EU-27's external-dependence structure in~2024: the concentration of
defence procurement suppliers and logistics counterparties varies
far more across member states than the concentration of financial,
energy, technology, or critical-input counterparties.  The
equal-weight arithmetic mean faithfully transmits this dispersion to
the composite.

Three consequences follow:

\begin{enumerate}[leftmargin=2em]
  \item \textbf{Policy leverage.}  Interventions that diversify
    defence procurement or logistics counterparties will have the
    largest effect on a country's composite rank.  Interventions in
    other domains will reduce the composite score but not the rank.

  \item \textbf{Measurement priority.}  Data quality in the defence
    and logistics axes is ranking-critical.  Measurement error in
    these domains propagates; measurement error in the financial axis
    does not.

  \item \textbf{Communication.}  The {\ISI} ranking is, in
    practice, a defence-and-logistics ranking with a residual
    contribution from the other four domains.  Users who cite the
    composite should disclose this concentration structure.
\end{enumerate}

\noindent
The finding is falsifiable in two ways.  First, re-estimation with
revised data (updated SIPRI transfers, expanded logistics coverage)
may alter the variance shares.  Second, application to a different
country set (e.g., OECD, G20) may produce a different variance
structure.  Until then, the duplex holds.
